\section{Introducción}

\subsection{Definición de microcuásar y analogía de escala con los cuásares}
\noindent Un microcuásar (término acuñado por I. F. Mirabel y L. F. Rodríguez en 1992) es un sistema estelar binario compuesto por dos cuerpos ligados gravitacionalmente: el remanente compacto de una estrella masiva (siendo esta una estrella de neutrones o un agujero negro de masa estelar) y una estrella compañera de la cual extrae materia. 

Tienen su origen en el descubrimiento de chorros de radio bipolares de escala de años luz proveniendo de la fuente de rayos X duros y $\gamma$ blandos 1E 1740.7-2942, conocida como el \textit{Gran Aniquilador} (que debe su nombre a la radiación de aniquilación electrón-positrón a $511$ keV en la región del centro galáctico \cite{mirabel,mornas2026bloque1}). Estos chorros, denominados \textit{jets}, y la radiación electromagnética de alta energía son resultado directo de la actividad de los microcuasares \cite{chaty,mirabel,paredes}. 

El estudio de este tipo de sistema binario es de gran utilidad para poder entender el comportamiento de otros como pueden ser las radiogalaxias y los cuásares extragalácticos, ya que se ha demostrado a través de observaciones que replican exactamente la morfología de lóbulos y chorros características de estos \cite{mirabel}. 

Dicha unidad física existente entre los cuásares, microcuásares y colapsares progenitores de estallidos de rayos $\gamma$ largos se debe a que comparten tres características fundamentales \cite{mirabel}:
\begin{itemize}
    \item[$\bullet$] Un objeto central compacto dotado de un pozo de potencial gravitatorio muy intenso y \textit{spin}.
    \item[$\bullet$] Un disco de acreción rotante calentado por disipación viscosa y turbulenta.
    \item[$\bullet$] \textit{Jets} impulsados por procesos magnetohidrodinámicos a lo largo del eje de rotación del sistema.
\end{itemize}

En el contexto de la relatividad general, los radios gravitacional $R_g$ y de Schwarzschild $R_S$ son los que fijan la escala espacial característica de un agujero negro \cite{chaty,mornas2026bloque1}:
\begin{equation}
    R_g=\frac{GM}{c^2}\simeq 1.48\,\left(\frac{M}{M_\odot}\right)\mbox{ km}\qquad R_S=2 R_g\simeq 2.95\,\left(\frac{M}{M_\odot}\right)\mbox{ km}
\end{equation}
Como además
\begin{equation}
    t_{din}\sim\frac{R_S}{c}=\frac{2GM}{c^3}
\end{equation}
se tiene que \cite{chaty,mirabel}
\begin{equation}
    \frac{t_{MC}}{t_{AGN}}=\frac{M_{MC}}{M_{AGN}}\sim 10^{-5}-10^{-8}
\end{equation}
Por tanto, los ciclos de acomplamiento acreción-eyección que en cuásares evolucionan en escalas temporales de años a siglos se pueden observar en microcuásares en escalas de meses a un año (ciclo completo de estallido) o incluso de segundos a minutos (en fenómenos rápidos). Además, los chorros de los cuásares alcanzan longitudes de kpc o Mpc ($\sim 10^{24}$ cm), mientras que en microcuásares se extienden típicamente hasta escalas de pc ($\sim 10^{17}$ cm). Este hecho da aún más peso a la importancia del estudio de microcuásares; gracias al paralelismo latente entre estas dos estructuras tan estrechamente relacionadas podemos observar en cuestión de meses lo que en cuásares tardaría años en evolucionar \cite{chaty}. 

Por otra parte, cuando el eje del \textit{jet} de un microcuásar apunta de forma prácticamente directa hacia la línea de visión terrestre, la amplificación Doppler relativista da lugar a un microblazar, análogo estelar de los blazares extragalácticos \cite{mirabel}.

\subsection{Anatomía de un sistema binario de rayos X y de un microcuásar}

\noindent Desde el punto de vista estructural, un microcuásar constituye un subconjunto de las binarias de rayos X en las que se desarrollan chorros relativistas emisores de radiación sincrotrón no térmica \cite{paredes}. Tal y como se detalla en \cite{paredes}, alrededor de un 22\% de las binarias de rayos X catalogadas clásicamente en la galaxia por Liu et al. (2006, 2007) han sido detectadas en ondas de radio. 

La anatomía completa de un microcuásar está compuesta por cinco elementos acoplados:
\begin{enumerate}
    \item \textbf{La estrella compañera} (también denominada \textbf{donante}). Atendiendo a su tipo espectral y su masa, podemos dividir los sistemas binarios en dos grandes familias:
    \begin{enumerate}
        \item \textbf{Binarias de Rayos X de Alta Masa (HMXB)}. Se trata de o bien una estrella joven, caliente y masiva de tipo espectral O o B ($M\geq 10 M_\odot$ y $T_{eff}\sim 10^4-5\times^4$ K) con intensa emisión ultravioleta; o bien de una estrella de Wolf-Rayet (tipo espectral W, correspondiente a núcleos de He y CNO con vientos estelares fuertes \cite{mornas2026bloque2}). 
        En este tipo de estrellas, la transferencia de masa hacia el objeto compacto se produce por la captura gravitatoria del viento estelar supersónico a través de un disco de decreción ecuatorial o por el desbordamiento del lóbulo de Roche.
        \item \textbf{Binarias de Rayos X de Baja Masa (LMXB)}. En este caso, la compañera es una estrella evolucionada o de secuencia principal de baja masa y de tipo espectral tardío ($M\leq 1-2 M_\odot$ y tipo posterior a B o gigante roja de baja masa \cite{fender_gallo,malzac}). 
        Como estas estrellas carecen de vientos intensos, la transferencia de masa es únicamente consecuencia del desbordamiento del lóbulo de Roche a través del punto interno de Lagrange $L_1$.
    \end{enumerate}
    \item \textbf{El objeto compacto acretor}. Responsable de la fuerza gravitatoria del sistema. Mediante medición por espectroscopía óptica o infrarroja del desplazamiento Doppler radial de las líneas de absorción de la estrella compañera se puede determinar su masa \cite{mornas2026bloque1}. Una vez calculada, se determinan las variables necesarias a partir de las cuales se obtiene la función de masa \cite{mornas2026bloque1} dada por las leyes de Kepler:
    \begin{equation}
        f(M_1)=\frac{(M_1\sin{i})^3}{(M_1+M_2)^2}=\frac{4\pi^2(a_2\sin{i})^3}{GP_{orb}^2}=\frac{P_{orb}K_2^3}{2\pi G}
    \end{equation}
    donde $M_1$ y $M_2$ son las masas del objeto compacto y de la estrella donante respectivamente, $i$ es la inclinación del plano orbital respecto al plano del cielo, $P_{orb}$ es el periodo orbital, $K_2$ es la semiamplitud de velocidad y $a_2\sin{i}$ es el semieje proyectado. 
    
    \newpage
    De esta forma, $f(M_1)$ constituye una cota inferior absoluta para $M_1$:
    \begin{enumerate}
        \item Si $M_1\leq M_{TOV}\simeq 2-3 M_\odot$ (estabilidad de Tolman-Oppenheimer-Volkoff \cite{mornas2026bloque2}) el objeto suele ser una estrella de neutrones, donde las fuentes de tipo Z emiten intensamente en radio de forma asociada a los \textit{jets}, las fuentes \textit{Atoll} lo hacen también más débilmente. Los púlsares de rayos X fuertemente magnetizados ($B\sim 10^{12}$ G) no producen \textit{jets} de radio debido al truncamiento a gran distancia que produce la magnetosfera sobre el disco de acreción \cite{chaty}.
        \item Si $M_1>3 M_\odot$ el objeto es un agujero negro de Kerr, caracterizado de forma única por su masa $M$ y su parámetro adimensional de \textit{spin} $a_\ast=\frac{Jc}{GM^2}\in[-1,1]$ al ser eléctricamente neutro \cite{mirabel,fender_gallo}. Esto se debe a que la presión de degeneración cuántica y las fuerzas nucleares fuertes son incapaces de detener el colapso.
    \end{enumerate} 
    \item \textbf{El disco de acreción frío}. Como el momento angular específico de la materia transferida desde la estrella compañera es muy elevado, el gas se asienta en un disco de acreción kepleriano en lugar de caer radialmente hacia el objeto compacto. Este disco es geométricamente delgado ($\frac{H}{r} \ll 1$) y ópticamente grueso ($\tau \gg 1$) \cite{malzac}. El momento angular es transportado hacia el exterior por las tensiones viscosas, que disipan energía potencial gravitaria provocando un calentamiento de las regiones internas del disco hasta $T\sim 10^6-10^7$ K. De esta forma, emite rayos X blandos ($\sim0,1-2$ keV) \cite{fender_gallo,malzac}.
    \item \textbf{La corona de plasma caliente} (flujo interno de acreción caliente). La emisión de cuerpo negro del disco es incapaz de explicar la emisión de fotones de rayos X duros y rayos $\gamma$ blandos que se detecta observacionalmente. Esto es debido a que la componente dura se origina en la corona, una región ópticamente delgada de electrones altamente energéticos ($T_e\sim10^9$ K $\approx100$ keV) \cite{malzac}.
    \item \textbf{\textit{Jets} y vientos del disco}. El sistema lanza \textit{jets} perpendicularmente al plano del disco de acreción, a lo largo del eje de rotación. Estos son observados en ondas de radio, milimétricas, infrarrojas y ópticas mediante radiación de sincrotrón y en rayos X y $\gamma$ mediante dispersión Compton inversa \cite{paredes,fender_gallo,malzac}. Recientemente, también se ha descubierto que el disco expulsa vientos ecuatoriales densos altamente ionizados no relativistas durante los estados espectrales blandos en los que el \textit{jet} polar se apaga \cite{fender_gallo}.
\end{enumerate}